\section{Totals: the capacity of a distance map is fixed by the number of bodies}\label{sec:totals}

\subsection{The distance map and its shape gradient}

Let $N\ge2$ point masses $m_i>0$ sit at $r_i\in\R^2$ with centre of mass at the origin and
\emph{unit inertia} $I:=\sum_i m_i\norm{r_i}^2=1$. Write $d_{ij}:=\norm{r_i-r_j}$, $M:=\sum_i m_i$,
and let $\ip{u}{v}_m:=\sum_i m_i\,u_i\cdot v_i$ be the mass metric on configurations. The
\emph{shape space} is the quotient of the unit-inertia, zero-centre-of-mass configurations by
translations (already removed), rotations and dilations; it carries the quotient Riemannian
metric, whose gradients we write $\nabla_{\mathrm{shape}}$. Distances are understood after the
inertia normalisation, so that $d_{ij}$ is a function on shapes.

\begin{lemma}[Inertia identity]\label{lem:inertia}
Without any centre-of-mass assumption,
$\sum_{i<j}m_im_j\norm{r_i-r_j}^2=M\,I-\bigl|\sum_im_ir_i\bigr|^2$; at zero centre of mass and
unit inertia, $\sum_{i<j}m_im_jd_{ij}^2=M$.
\T\ (expand the square and complete it; the two correction sums cancel).
\end{lemma}

\begin{theorem}[Per-pair capacity]\label{thm:perpairs}
Assume no two bodies coincide. Then
\begin{equation}\label{eq:perpairs}
   \norm{\nabla_{\mathrm{shape}}d_{ij}}_m^2=\frac{1}{m_i}+\frac{1}{m_j}-d_{ij}^2 .
\end{equation}
\T\ Proof below; independent numerical verification \V\ gate \gate{g01\_shape\_spectrum} (g),
five configurations with unequal masses up to $N=8$, maximum residual $8.9\times10^{-16}$.
\end{theorem}

\begin{proof}
Let $G^{ij}$ be the mass-metric gradient of $d_{ij}$ on configuration space. Its only nonzero
particle components are $(r_i-r_j)/(m_id_{ij})$ and $-(r_i-r_j)/(m_jd_{ij})$, so
$\norm{G^{ij}}_m^2=1/m_i+1/m_j$. Three identities follow by direct evaluation:
$\langle G^{ij},\text{translations}\rangle_m=0$ (so $G$ is horizontal for translations),
$\ip{G^{ij}}{\rho r}_m=d_{ij}\,e_{ij}\cdot Je_{ij}=0$ with $\rho r$ the rotation field and
$e_{ij}=(r_i-r_j)/d_{ij}$ (rotation invariance), and $\ip{G^{ij}}{r}_m=e_{ij}\cdot(r_i-r_j)=d_{ij}$
(Euler's identity for the $1$-homogeneous function $d_{ij}$). Since $I=1$, the dilation direction
$r$ is a unit vector, and $H^{ij}:=G^{ij}-d_{ij}r$ is therefore horizontal and orthogonal to
translations, rotation and dilation. On a horizontal variation the unit-inertia constraint is
stationary, so $H^{ij}$ is the shape gradient and
$\norm{H^{ij}}_m^2=\norm{G^{ij}}_m^2-2d_{ij}\ip{G^{ij}}{r}_m+d_{ij}^2\norm r_m^2=1/m_i+1/m_j-d_{ij}^2$,
using $\norm r_m^2=I=1$.
\end{proof}

\begin{theorem}[Capacity sum rule]\label{thm:sumrule}
Under the hypotheses of Theorem~\ref{thm:perpairs},
\begin{equation}\label{eq:sumrule}
   \sum_{i<j}\norm{\nabla_{\mathrm{shape}}d_{ij}}_m^2
   =(N-1)\sum_i\frac1{m_i}-\sum_{i<j}d_{ij}^2 ,
\end{equation}
and for $N$ equal unit masses this is $N(N-1)-N=N(N-2)$.
\T\ (sum \eqref{eq:perpairs} over the $\binom N2$ pairs; each index occurs $N-1$ times, and
Lemma~\ref{lem:inertia} gives $\sum_{i<j}d_{ij}^2=N$). \V\ Gate \gate{g01\_shape\_spectrum} (f),
$N=3,\dots,7$, residuals below $2\times10^{-6}$ (finite-difference floor).
\end{theorem}

\begin{corollary}[The total is configuration-free]\label{cor:total}
For fixed masses and unit inertia the right-hand side of \eqref{eq:sumrule} does not depend on
the configuration at all. Equivalently, the \emph{mean} squared singular value of the distance
map over the $(2N-4)$-dimensional shape space is $N/2$ for equal unit masses, whatever the shape.
\T
\end{corollary}

\begin{remark}[What the sum rule is and is not]\label{rem:sumrule-meaning}
It is a conservation law, not a capacity formula: it says the total squared response of all
pairwise distances to a shape deformation is fixed by the number of bodies, and the shape only
decides how that total is distributed among the $\binom N2$ distances. The corpus's
\emph{per-shape} capacity claims --- a plateau of $13$ against an aperture prediction of $12.57$,
a cell value $8$ against $8.17$, a fitted coefficient $C_0$ absorbing a tube-counting
normalisation --- are statements about the \emph{distribution} and therefore necessarily carry a
fitted constant. \V/\Hh\ The conservation law is what survives without one. This is the sense in
which the corpus's most quotable idea (``a transport can forget almost every direction and still
retain a constant'') is true: the constant is the total, and it is combinatorial.
\end{remark}
